% Generated draft: requires mathematical and source review.
\documentclass[11pt]{article}
\usepackage[margin=25mm]{geometry}
\usepackage{fontspec}
\setmainfont{Latin Modern Roman}
\usepackage{amsmath,amssymb,mathtools,unicode-math}
\setmathfont{Latin Modern Math}
\usepackage{enumitem,needspace,xurl,hyperref}
\setlength{\parindent}{0pt}
\setlength{\parskip}{5pt}
\setlength{\emergencystretch}{3em}
\newcommand{\N}{\mathbb{N}}
\newcommand{\Z}{\mathbb{Z}}
\newcommand{\Q}{\mathbb{Q}}
\newcommand{\R}{\mathbb{R}}
\newcommand{\C}{\mathbb{C}}
\newcommand{\F}{\mathbb{F}}
\newcommand{\A}{\mathbb{A}}
\newcommand{\PP}{\mathbb{P}}
\newcommand{\E}{\mathbb{E}}
\newcommand{\eps}{\varepsilon}
\newcommand{\dd}{\,\mathrm{d}}
\begin{document}
\section*{UnsolvedMath ID: 30004587}
\textbf{Source problem code:} catalogue entry
\subsection{Definitions and mathematical statement}
A quantum Markov semigroup is a family \(P=(P_t)_{t\geq 0}\) of unital completely positive maps on an operator algebra, satisfying \(P_0=\mathrm{id}\) and \(P_{s+t}=P_sP_t\). The relevant algebra class, continuity requirements, and any normality, invariant-state, traciality, or symmetry assumptions must be specified to determine the intended setting; the supplied record does not specify them.

For two such semigroups \(P^{(1)}\) and \(P^{(2)}\), their tensor-product semigroup is formally given by
\[
(P^{(1)}\otimes P^{(2)})_t(a\otimes b)
   =P^{(1)}_t(a)\otimes P^{(2)}_t(b),\qquad t\geq0.
\]
The operator-algebra tensor product and the extension of these maps to it are part of the unspecified setting.

Write \(P\models\operatorname{GE}(K,\infty)\) for satisfaction of the ordinary quantum Markov gradient estimate denoted by this notation in the cited report. Its defining inequality, gradient structure, auxiliary choices, and domains are absent from the supplied record, so this notation remains an unresolved predicate here.

\textbf{Source question, in schematic form.} Let \(\mathcal C\) denote the intended class of quantum Markov semigroups and \(\mathcal K\) the intended range of curvature parameters. Is it true that, for every \(K\in\mathcal K\) and every \(P^{(1)},P^{(2)}\in\mathcal C\),
\[
\left[
 P^{(1)}\models\operatorname{GE}(K,\infty)
 \ \text{and}\ 
 P^{(2)}\models\operatorname{GE}(K,\infty)
\right]
\Longrightarrow
P^{(1)}\otimes P^{(2)}\models\operatorname{GE}(K,\infty)?
\]
This displays the common-\(K\), two-factor interpretation of tensor stability. The record does not explicitly specify this quantifier convention, the classes \(\mathcal C\) and \(\mathcal K\), or the definition of \(\operatorname{GE}\); consequently it does not determine a fully self-contained mathematical proposition. The target is the ordinary estimate \(\operatorname{GE}\). The supplied literature assessment distinguishes it from the stronger complete estimate \(\operatorname{CGE}\).
\subsection{Short English statement}
The question asks whether tensoring quantum Markov semigroups preserves the ordinary gradient estimate \(\operatorname{GE}(K,\infty)\), with the same curvature parameter under the displayed interpretation. The supplied record omits the estimate's definition and essential hypotheses, so the statement remains schematic and requires mathematical review. Its open status is supplied dated curation, not an independently verified conclusion.
\subsection{Sources}
\begin{itemize}
\item Supplied local UnsolvedMath registry record, permanent ID \(30004587\), title ``Tensor Stability of Quantum Markov Gradient Estimates,'' exact source problem code \texttt{catalogue entry}; collection: Oberwolfach Reports Open Problems. No individual proposer is supplied. The record gives proposed year \(2021\).
\item Supplied source citation: \emph{Geometry, Dynamics and Spectrum of Operators on Discrete Spaces} (2022), Oberwolfach Reports, pp.\ 33--85, DOI \texttt{10.4171/owr/2021/2}. \url{https://doi.org/10.4171/owr/2021/2}. The supplied background describes these reports as documenting mathematical workshops and their open-problem discussions.
\item Supplied literature reference: Wirth and Zhang, ``complete estimate'' (the record's label), DOI \texttt{10.1007/s00220-021-04199-4}. \url{https://doi.org/10.1007/s00220-021-04199-4}. No full bibliographic title is supplied.
\item Literature assessment checked at \texttt{2026-08-22}: supplied dated curation, not independently verified. It labels the ordinary \(\operatorname{GE}(K,\infty)\) tensor-stability question open and states that the cited Wirth--Zhang result proves tensor stability for the stronger complete estimate \(\operatorname{CGE}\), rather than the ordinary estimate asked about here. The supplied background explicitly requests independent verification of time-sensitive status and cited claims.
\end{itemize}
\end{document}
